% from ega4-2-fr.tex
\oldpage[IV-2]{220}

en particulier si $X$ est \emph{localement immersible dans un préschéma
régulier} (5.6.4)), $X$ est un \emph{préschéma excellent} (7.8.5).
\end{proposition}

\begin{proof}
Cela résulte aussitôt de (7.9.4) et (7.9.2).
\end{proof}

\begin{remark}[7.9.6]
\label{IV.2.s7.9.6-fr}
Il est possible que la réciproque de (7.9.5) soit vraie, c'est-à-dire que
l'hypothèse que $X$ est réduit et que les anneaux des ouverts affines de $X$
vérifient les conditions (ii) et (iii) de (7.8.2) implique la possibilité de
résoudre $X$ (et par suite aussi la possibilité de résoudre tout préschéma
réduit localement de type fini sur $X$, en vertu de (7.4.4) et (6.12.4)).
C'est vrai en tout cas lorsqu'on se borne aux préschémas noethériens réduits
dont les corps résiduels sont de \emph{caractéristique} $0$, comme le montrent
les résultats récents de Hironaka [35] (ce dernier énonce ses résultats sous
des hypothèses trop restrictives, ses raisonnements étant en fait valables
sous les seules hypothèses (ii) et (iii) de (7.8.2) pour les anneaux des
ouverts affines de $X$, jointes au fait que les corps résiduels sont de
caractéristique $0$). Au contraire, pour le cas général, on ne sait jusqu'ici
résoudre que les schémas algébriques de dimension $2$ sur un corps parfait
[36]. Vu l'importance que la résolution des singularités est en train de
prendre dans l'étude topologique des schémas (notamment en ce qui concerne
leurs propriétés homologiques et homotopiques), cela donne un intérêt
particulier à la catégorie des anneaux excellents ou des préschémas
excellents.

On peut aussi noter à ce propos que la partie la moins délicate des
raisonnements de Hironaka (\emph{loc, cit.}) montre qu'indépendamment de toute
hypothèse sur les caractéristiques des corps résiduels, et en utilisant
uniquement les propriétés (ii) et (iii) de (7.8.2) pour les anneaux locaux
d'un préschéma $X$, le problème de la résolution des singularités de $X$ se
ramène au cas où $X=\operatorname{Spec}(A)$, $A$ étant un anneau local
noethérien intègre et \emph{complet}. Par conséquent, si des résultats
ultérieurs devaient mettre en défaut la conjecture avancée plus haut, et
devaient donc amener à formuler des conditions plus restrictives sur les
anneaux locaux de $X$, cela ne pourrait guère être que par l'intermédiaire de
conditions restrictives imposées à des anneaux locaux complets (probablement
des conditions concernant leurs corps résiduels, peut-être la condition
$[k:k^p]<+\infty$ ($p$ exposant caractéristique de $k$)).
\end{remark}

\begin{env}[7.9.7]
\label{IV.2.s7.9.7-fr}
Considérons d'une part une sous-catégorie pleine $\mathcal C$ de la catégorie
des préschémas localement noethériens, d'autre part une propriété
$\boldsymbol R(A)$, soumises aux conditions suivantes :

$1^\circ$ Pour tout $X\in\mathcal C$, tout préschéma localement de type fini
sur $X$ appartient à $\mathcal C$. Pour tout anneau noethérien $A$ tel que
$\operatorname{Spec}(A)\in\mathcal C$ et toute partie multiplicative $S$ de
$A$, on a $\operatorname{Spec}(S^{-1}A)\in\mathcal C$.

$2^\circ$ Pour tout $X\in\mathcal C$, l'ensemble
$U_{\boldsymbol R}(X)$ des $x\in X$ tels que
$\boldsymbol R(\mathcal O_{X,x})$ soit vraie est ouvert dans $X$.

$3^\circ$ Pour tout anneau local noethérien $A$ tel que
$\operatorname{Spec}(A)\in\mathcal C$ et tout élément régulier $t$ de l'idéal
maximal de $A$, $\boldsymbol R(A/tA)$ entraîne $\boldsymbol R(A)$.

Désignons alors comme dans (7.5.0) par $\boldsymbol P(Z,k)$, pour un corps $k$
et un $k$-préschéma $Z\in\mathcal C$, la propriété suivante :

« Pour toute extension finie $k'$ de $k$, tous les anneaux locaux de
$Z\otimes_k k'$ vérifient la propriété $\boldsymbol R$. »

